\chapter{Conclusiones}
\label{ch:conclusions}

\section{Conclusiones parciales}

El proyecto Litoteca demuestra la viabilidad de integrar un modelo de clasificación
de rocas en una aplicación web orientada a usuarios no expertos. Hasta el estado
actual, se completaron los componentes estructurales más importantes: definición del
contexto, selección y limpieza del dataset, construcción de la interfaz, proxy de
predicción, despliegue en Cloudflare Pages y preparación del entrenamiento en Custom
Vision.

La preparación de datos fue un hallazgo crítico. El export original contenía
aumentos sintéticos y fuga entre splits; sin la deduplicación por imagen original,
las métricas habrían sido optimistas y poco defendibles. La versión limpia generada
para entrenamiento reduce este riesgo y deja una base más sólida para evaluar el
modelo.

La aplicación cumple el criterio de utilidad práctica: no solo intenta clasificar,
sino que explica el resultado con fichas, confianza, alternativas y contexto
geológico. Esto es especialmente importante en un dominio donde clases visualmente
similares pueden confundirse.

\section{Limitaciones}

\begin{itemize}
  \item La iteración Quick disponible tiene recall global moderado (52.43\%), con
        clases débiles como shale, clay y diatomite; requiere mejora antes de
        considerarla robusta.
  \item Aún faltan métricas complementarias calculadas sobre test (F1-score, ROC/AUC
        y matriz de confusión) además de las métricas nativas de Custom Vision.
  \item El dataset cubre solo 9 clases y no incluye rocas metamórficas.
  \item Existen clases síliceas con alta similitud visual y conflictos históricos de
        etiquetado en la fuente.
  \item La predicción en producción depende de la disponibilidad del endpoint de
        Custom Vision y de la configuración de secretos en Cloudflare Pages.
\end{itemize}

\section{Trabajo pendiente}

Para la entrega/demo inmediata y la mejora posterior:

\begin{enumerate}
  \item Publicar o conectar la iteración Quick actual en Cloudflare Pages para la
        demo del día.
  \item Ejecutar validación sintética por API con \texttt{data/prepared-clean-train/test}
        y registrar aciertos, errores y confianza.
  \item Realizar prueba de campo controlada para el video/demo.
  \item Lanzar una iteración adicional (Advanced o dataset ajustado) para mejorar
        recall en clases débiles; si mejora, reemplazar la iteración conectada.
  \item Actualizar este informe con F1-score, ROC/AUC, matriz de confusión y
        conclusiones finales.
\end{enumerate}
